% TOP_PROBLEM: {"id": 30006759, "title": "NP versus the Existential Theory of the Reals", "category_id": 15, "category": {"id": 15, "name": "computer_science", "display_name": "Computer Science", "description": "Computational complexity, algorithms, and theoretical CS.", "slug": "computer-science", "order_index": 15, "created_at": "2026-07-31T15:26:25.671Z"}, "status": "open"}
% FIRST_POSTED: 2026-09-27
% !TeX program = lualatex
% Compile twice: lualatex 30006759.tex
% All references and prose are embedded; no BibTeX database or external inputs.
\documentclass[11pt,a4paper]{article}
\usepackage[margin=24mm,headheight=14pt]{geometry}
\usepackage{fontspec}
\setmainfont{Latin Modern Roman}
\setsansfont{Latin Modern Sans}
\setmonofont{Latin Modern Mono}
\usepackage{amsmath,amssymb,mathtools}
\usepackage{unicode-math}
\setmathfont{Latin Modern Math}
\usepackage{microtype}
\usepackage{enumitem,xurl,needspace}
\usepackage[dvipsnames]{xcolor}
\usepackage{hyperref}
\hypersetup{unicode=true,colorlinks=true,linkcolor=MidnightBlue,urlcolor=MidnightBlue,
 pdftitle={500 Mathematical Problem Statements},pdfauthor={Source-annotated compilation},bookmarksnumbered=true}
\usepackage{bookmark}
\usepackage{tocloft}
\setlength{\cftsecnumwidth}{3em}
\cftsetpnumwidth{2.5em}
\cftsetrmarg{4em}
\usepackage{titlesec}
\titleformat{\section}{\Large\bfseries\raggedright}{\thesection}{0.7em}{}
\usepackage{fancyhdr}
\pagestyle{fancy}\fancyhf{}
\fancyhead[L]{\small\sffamily Mathematical problem statements}
\fancyhead[R]{\small Source edition 22}
\fancyfoot[C]{\thepage}
\setlength{\parindent}{0pt}
\setlength{\parskip}{5pt plus 1pt}
\setlength{\emergencystretch}{3em}
\setcounter{tocdepth}{1}
\setcounter{secnumdepth}{1}
\setlist[itemize]{leftmargin=3em,itemsep=5pt,topsep=4pt,parsep=0pt}
\allowdisplaybreaks
\newcommand{\N}{\mathbb{N}}
\newcommand{\Z}{\mathbb{Z}}
\newcommand{\Q}{\mathbb{Q}}
\newcommand{\R}{\mathbb{R}}
\newcommand{\C}{\mathbb{C}}
\newcommand{\F}{\mathbb{F}}
\newcommand{\A}{\mathbb{A}}
\newcommand{\PP}{\mathbb P}
\newcommand{\E}{\mathbb{E}}
\newcommand{\eps}{\varepsilon}
\newcommand{\dd}{\,\mathrm d}
\newcommand{\sref}[2]{\hyperlink{source:#1:#2}{[S#2]}}
\newcommand{\eref}[2]{\hyperlink{extra:#1:#2}{[E#2]}}
% Turn the next switch off for a shorter reading copy. The quotations preserve
% every exactTarget field, including qualifications omitted from a short summary.
\newif\ifshowcatalogue
\showcataloguetrue
\newcommand{\cataloguescope}[1]{\ifshowcatalogue
 \paragraph{Exact scope in the supplied catalogue.}
 \begin{quote}\small #1\end{quote}\fi}

% Standalone research-notebook wrapper; the source preamble above is retained.
\numberwithin{equation}{section}
\hypersetup{pdftitle={Research attempt -- problem 30006759},
 pdfauthor={Research notebook; original compilation retained}}
\fancyhead[L]{\small\sffamily Research notebook}
\fancyhead[R]{\small\sffamily Problem 30006759 | 27 September 2026}
\begin{document}
\setcounter{section}{0}
% ================================================================
% problemId: 30006759
% Canonical problem ID: 30006759
% exactTarget SHA-256: bae376ace0c50f80edf1a296025017c06149c0e0bf93a3d465a5aec813b81349
\clearpage
\section*{\texorpdfstring{NP versus the Existential Theory of the Reals}{NP versus the Existential Theory of the Reals}}\label{30006759}
\subsection{Definitions and mathematical statement}
The existential theory of the reals asks whether a sentence
$\exists x_1,\ldots,x_m\in\R\;\Phi(x_1,\ldots,x_m)$ is true, where $\Phi$ is a quantifier-free Boolean combination of polynomial equalities and inequalities with integer coefficients given in finite binary encoding. The complexity class $\exists\R$ consists of languages polynomial-time many-one reducible to this decision problem. With $\mathrm{NP}$ defined by polynomial-length binary certificates and polynomial-time deterministic verification, the selected question is
\[
 \exists\R\subseteq\mathrm{NP},\qquad\text{equivalently }\mathrm{NP}=\exists\R.
\]
Witnesses for the real formula itself need not be rational or have short numerical encodings; the question concerns existence of \emph{some} NP verification scheme, not merely writing down a real solution.
\par\smallskip\textit{Formulation and concept references:} \sref{30006759}{1}.
\cataloguescope{Determine whether every problem decidable in the existential theory of the reals (exists R) can be decided in nondeterministic polynomial time, i.e., decide whether NP = exists R.}
\subsection{Short English statement}
Can the existence of solutions to finite systems of real polynomial constraints always be certified by a short binary proof that is efficiently checkable?
% BEGIN RESEARCH ADDITION ATTEMPT-158
\subsection{Research attempt}

\paragraph{Current frontier.}
The established inclusions are
$\mathrm{NP}\subseteq\exists\R\subseteq\mathrm{PSPACE}$. The compendium
emphasizes both the absence of a general short-certificate theorem and the
fact that large coordinate descriptions do not themselves separate these
classes \sref{30006759}{1}, Sections 1.2--1.3. A useful adjacent problem is
\textsf{PosSLP}: deciding positivity of an integer represented by a
straight-line arithmetic program. Allender--B\"urgisser--Kjeldgaard-Pedersen--
Miltersen place it in the counting hierarchy, but that is not an NP bound
\eref{30006759}{1}. No verified 2024--2026 result located for this entry settles
the selected equality. The attempt below constructs a genuine binary
certificate for a quantitative subclass and tests the proposed extensions.

\paragraph{Attack route.}
Literal rational solutions fail when only irrational solutions exist, and
even strictly feasible rational solutions can require exponentially many
bits. Algebraic towers can be short, but efficient general sign checking
is an additional problem. A third route is to certify a solution
\emph{indirectly} by a rational contraction box. Its prospective advantage
is that no exact encoding of the solution is needed. Its missing input is
a polynomial-size covering theorem: every satisfiable instance would need
a short certificate of this kind, or a controlled generalization covering
singular and poorly conditioned solutions.

\paragraph{Attempt: reducing representation issues without hiding cost.}
Introducing one variable for each arithmetic gate replaces polynomial
computations by quadratic gate equations. Binary-encoded exponents can be
evaluated by repeated squaring with a number of gates polynomial in their
bit lengths. A guessed consistent truth assignment to the Boolean
combination selects sign conditions on the computed outputs. This explains
why quadratic systems are a useful test class; it does not manufacture
short values for the added real variables. In particular, the computation
model throughout this entry is ordinary binary computation, not unit-cost
arithmetic on arbitrary real or exponentially long integer registers.

For the positive lemma, consider the explicit subclass
\begin{equation}\label{eq158-system}
 P_1(x)=\cdots=P_m(x)=0,\qquad Q_j(x)>0\quad(1\le j\le s),
 \qquad x\in\R^m,
\end{equation}
where all $P_i,Q_j$ are rational polynomials of degree at most two, given
by their coefficients. Clearing positive denominators gives an instance
of the target problem. The equal number of equations and variables is a
hypothesis for this lemma, not a normal form asserted for every instance.

\paragraph{Attempt: a rational geometric certificate.}
A certificate consists of $x_0\in\Q^m$, $r\in\Q_{>0}$, and an invertible
matrix $A\in M_m(\Q)$. Let
\[
 B=\{x:\|x-x_0\|_\infty\le r\},\qquad P=(P_1,\ldots,P_m)^t,
 \qquad H_k=\partial_kDP.
\]
The matrices $H_k$ are constant because $P$ is quadratic. All matrix norms
in the next display are induced $\ell^\infty$ norms. Compute the rational
numbers
\begin{align}
 q&=\|I-A DP(x_0)\|_\infty+
       r\sum_{k=1}^m\|A H_k\|_\infty,\label{eq158-q}\\
 M_j&=\|DQ_j(x_0)\|_1+
       r\sum_{k=1}^m\|\partial_k DQ_j\|_1.\label{eq158-M}
\end{align}
The verifier checks, by exact rational arithmetic,
\begin{equation}\label{eq158-checks}
 \det A\ne0,\quad q\le\tfrac12,\quad
 \|A P(x_0)\|_\infty\le r/2,\quad Q_j(x_0)-rM_j>0\ (1\le j\le s).
\end{equation}

\textbf{Certificate lemma.} Every accepted certificate proves that
\eqref{eq158-system} is satisfiable. Verification takes time polynomial in
the combined binary length of the instance and certificate.

\textit{Proof of soundness.} Put $T(x)=x-A P(x)$. For $x\in B$,
\[
 DP(x)=DP(x_0)+\sum_k(x_k-x_{0,k})H_k,
 \qquad\|DT(x)\|_\infty\le q.
\]
The segment between any two points of $B$ lies in $B$, so the mean value
formula proves that $T$ is $q$-Lipschitz on $B$. Also
\[
 \|T(x)-x_0\|_\infty
 \le\|AP(x_0)\|_\infty+q\|x-x_0\|_\infty\le r.
\]
Thus $T$ maps the complete closed box to itself and is a contraction.
Successive iteration is Cauchy because successive differences decrease
geometrically; its limit $x_*$ satisfies $T(x_*)=x_*$. Since $A$ is
invertible, $P(x_*)=0$. Equation~\eqref{eq158-M} bounds
$\sup_{x\in B}\|DQ_j(x)\|_1$, so
\[
 Q_j(x_*)\ge Q_j(x_0)-rM_j>0.
\]
This proves all the required conditions. There is no need to compute the
limit, approximate its coordinates to an unbounded precision, or test a
real number for exact equality.

\textit{Proof of verification complexity.} Polynomial evaluation at
rational points of supplied bit length, rational matrix multiplication,
row sums of absolute values, and rational comparisons require polynomially
many integer operations on polynomial-bit integers here. Exact determinant
computation also has polynomial bit complexity, for example by
fraction-free elimination. The finite sums contain only polynomially many
terms. Hence the verification bound is polynomial in the \emph{supplied}
encoding length. This statement alone does not bound that length by a
polynomial in the original instance. \hfill$\square$

An irrational solution is not an obstacle to this certificate. For
$P(x)=x^2-2$ and $Q(x)=x$, take
\[
 x_0=\tfrac32,\qquad r=\tfrac14,\qquad A=\tfrac13.
\]
Then $q=1/6$, $|AP(x_0)|=1/12<r/2$, and
$Q(x_0)-rM=5/4>0$. The certificate proves existence of the positive root
without encoding $\sqrt2$ as a binary number. These inequalities are also
checked in the accompanying exact-arithmetic script.

\paragraph{Attempt: when polynomial-length boxes do follow.}
Here is a quantitative subclass for which the preceding verifier really
does give short certificates. Let $N$ be the instance length and suppose
there is a feasible $x_*$ with invertible $DP(x_*)$, and a fixed polynomial
$h$ such that
\begin{equation}\label{eq158-condition}
 \|x_*\|_\infty,\ \|DP(x_*)^{-1}\|_\infty\le2^{h(N)},
 \qquad \min_j Q_j(x_*)\ge2^{-h(N)}.
\end{equation}
For $s=0$ omit the minimum condition. The quadratic coefficients and the
number of variables have bounds of the same single-exponential type from
the input encoding. On a unit neighborhood of $x_*$ all first and second
derivatives are therefore bounded by $2^{g(N)}$ for a fixed polynomial
$g$, after increasing $g$ as necessary.

Choose a dyadic radius $r=2^{-k(N)}$ sufficiently small that the Hessian
contribution in \eqref{eq158-q}, evaluated with
$A_*=DP(x_*)^{-1}$, is at most $1/8$, and that all inequalities retain a
margin on the radius-$2r$ box. Approximate $x_*$ and $A_*$ by dyadic
$x_0,A$ to an accuracy $2^{-l(N)}$, with $l$ another sufficiently large
fixed polynomial. The derivative bounds and the bounds on $A_*$ imply
\[
 \|I-A DP(x_0)\|_\infty\le\tfrac18,\qquad
 \|A P(x_0)\|_\infty\le r/2,
\]
and keep the inequalities in \eqref{eq158-checks} strict. For example,
choose the rounding errors smaller than $r$ times the reciprocal of a
suitable product of the finitely many bounds $2^{g(N)}$. Such a choice
still has polynomial bit length. Finally $\|I-A DP(x_0)\|<1$ implies that
$A DP(x_0)$, and hence the square matrix $A$, is invertible. This constructs
a polynomial-length certificate under \eqref{eq158-condition}.

No claim is made that every feasible system has such a point. Singular
solutions, positive-dimensional solution sets, or very small inequality
margins require another argument. Converting them to the square regular
case by slicing or deflation without uncontrolled growth is itself a
missing lemma.

\paragraph{Stress test: small coordinates even in a bounded open set.}
Consider the strictly feasible quadratic system
\begin{equation}\label{eq158-small}
 0<x_0<\tfrac12,\qquad 0<x_{i+1}<x_i^2\quad(0\le i<n).
\end{equation}
All coordinates lie in $(0,1)$, but induction gives
$0<x_n<2^{-2^n}$. If $x_n=a/b$ is a reduced positive rational, then
$a\ge1$ and $b>2^{2^n}$, so its ordinary denominator encoding has more
than $2^n$ bits. Boundedness and strict feasibility do not imply a
polynomial-length literal rational witness.

Nevertheless this family has a short \emph{description} of a witness:
$x_0=1/4$ and $x_{i+1}=x_i^2/2$, equivalently
$x_i=2^{-(3\cdot2^i-1)}$. The recurrence proves all constraints.
Therefore \eqref{eq158-small} is not a complexity separation; it only
refutes a particular coordinate-encoding proposal.

A different obstruction affects a single primitive algebraic-number
encoding. The system
\[
 y_0^2=2,\quad y_0>0,\qquad y_{i+1}^2=y_i,\quad y_{i+1}>0
       \quad(0\le i<n)
\]
has all coordinates in $(1,2)$, yet $y_n$ has degree $2^{n+1}$ over $\Q$.
Indeed its minimal polynomial is $T^{2^{n+1}}-2$, irreducible by
Eisenstein at $2$. A tower description is short. In fact such systems may
also be amenable to geometric certificates; exponential algebraic degree
is not a proof that all certificate systems are long.

\paragraph{Stress test: the price of arithmetic compression.}
Let an integer straight-line program start from $1$ and use addition,
subtraction, and multiplication. Introduce one real variable per gate and
impose its defining equation, such as $z_k-z_i z_j=0$, together with the
initial constant equation and $z_{\mathrm{out}}>0$. Induction on the gates
shows that this polynomial-size existential formula is satisfiable exactly
when the program's integer output is positive. Thus
\begin{equation}\label{eq158-posslp}
 \mathrm{NP}=\exists\R\quad\Longrightarrow\quad
 \textsf{PosSLP}\in\mathrm{NP}.
\end{equation}
For integer output $z$, the complement condition $z\le0$ is equivalent to
$1-z>0$, obtained by two extra gates. Hence the same implication would
actually put $\textsf{PosSLP}$ in $\mathrm{NP}\cap\mathrm{coNP}$.
The reduction is proved here; the role of this sign problem in numerical
complexity is established in \eref{30006759}{1} and discussed in
\sref{30006759}{1}.

Consequently a proposed universal certificate consisting of succinct
arithmetic programs cannot simply declare their signs efficiently
checkable. Exact cancellation may involve integers with exponentially
many bits. Such a declaration assumes a major part of the problem rather
than solving it. Conversely, the two explicit large-description families
above do not show $\textsf{PosSLP}\notin\mathrm{NP}$ or
$\mathrm{NP}\ne\exists\R$. Nor does \eqref{eq158-posslp} imply
$\mathrm P=\mathrm{NP}$; the latter is a different, stronger algorithmic
assertion. Positivstellensatz certificates of \emph{infeasibility} address
the opposite side and do not by themselves provide the required
satisfiability certificates.

\paragraph{Outcome.}
\textbf{Outcome: Exploratory approach with unresolved gap.}

The attempt supplies a sound polynomial-time binary verifier for rational
contraction certificates, proves polynomial certificate length under
explicit regularity and margin bounds, and gives exact small examples.
It also proves why two naive witness-size arguments fail, and identifies
a precise compressed-sign consequence of the full target. No novelty is
claimed for the underlying contraction principle or the standard
straight-line-program reduction.

What remains missing is a universal polynomial-size certificate theorem,
including singular feasible sets and arbitrarily ill-conditioned or
narrow solution regions, without assuming polynomial-time succinct sign
arithmetic. The large-coordinate and large-degree constructions only rule
out specified formats; they do not rule out NP verification in general.
Thus neither equality nor separation of the two classes has been proved.
% END RESEARCH ADDITION ATTEMPT-158
\subsection{Sources}
\begin{itemize}\raggedright
\item[\textnormal{[S1]}]\hypertarget{source:30006759:1}{} M. Schaefer, J. Cardinal, T. Miltzow, The Existential Theory of the Reals as a Complexity Class: A Compendium, arXiv:2407.18006, 2024.
\url{https://arxiv.org/abs/2407.18006}.
{\small\textit{Catalogue formulation source.}}
% BEGIN RESEARCH ADDITION SOURCES-158
\item[\textnormal{[E1]}]\hypertarget{extra:30006759:1}{} Eric Allender, Peter
B\"urgisser, Johan Kjeldgaard-Pedersen, and Peter Bro Miltersen,
\emph{On the Complexity of Numerical Analysis}, \emph{SIAM Journal on
Computing} 38(5) (2009), 1987--2006, Introduction and the PosSLP
counting-hierarchy theorem.
\url{https://doi.org/10.1137/070697926}.
Author manuscript: \url{https://www.cs.rutgers.edu/~allender/papers/slp.pdf}.
{\small\textit{Consulted 27 September 2026; a counting-hierarchy bound is
not treated as membership in NP.}}
% END RESEARCH ADDITION SOURCES-158
\end{itemize}

\end{document}
